% Source: https://artofproblemsolving.com/community/c3h2300946_theory_of_equations
% Stable identity: aops-topic-2300946-q01
If $x=0$, the first equation gives $y=0$; if $y=0$, the second gives $x=0$. So $(0,0)$ is a solution, and otherwise $x\ne0$ and $y\ne0$.

Now let $s=x+y$ and multiply the two equations:
\[3x(s-2)\cdot y(s-1)=2y\cdot9x.\]
Dividing by $3xy\ne0$, $(s-2)(s-1)=6$, i.e. $s^2-3s-4=0$, so $s=4$ or $s=-1$.
\begin{itemize}
\item $s=4$: the first equation gives $6x=2y$, so $y=3x$, and $x+3x=4$ gives $(x,y)=(1,3)$.
\item $s=-1$: the first equation gives $-9x=2y$, so $y=-\dfrac92x$, and $x-\dfrac92x=-1$ gives $x=\dfrac27$, $y=-\dfrac97$.
\end{itemize}
Both satisfy the second equation as well: $3\cdot(4-1)=9=9\cdot1$, and $-\dfrac97\cdot(-2)=\dfrac{18}{7}=9\cdot\dfrac27$. So the solutions are
\[\boxed{(x,y)=(0,0),\ (1,3),\ \left(\dfrac27,-\dfrac97\right)}.\]
